A pulsar located 1000\,pc far from Earth, 10\,000 times more luminous than our
Sun, emits radiation only from its two opposite poles, creating an homogeneous
emission beam shaped as double cone with opening angle $\alpha = 4^\circ$.
Assuming the angle between the rotation axis and the emission axis is
$30^\circ$, and assuming a random orientation of the pulsar beams in relation
to an observer on Earth, what is the probability of detecting the pulses? In
case we can see it, what is the apparent bolometric magnitude of the pulsar?
